\section{Out-of-Fold Label Filtering}\label{sec:method}

\paragraph{Setting.}
The filter changes which samples logistic regression is fit on. The training set $D=\{(x_i,\tilde{y}_i)\}_{i=1}^{n}$ has $n$ samples, each a feature vector $x_i$ with a given label $\tilde{y}_i$ that may be wrong. The base model $f_D$ maps a feature vector to a class. It chains standardization and logistic regression, both fit on the set in its subscript. The logistic regression baseline is $f_D$, trained on all given labels without filtering.

\paragraph{Out-of-fold check.}
A model that has never seen a sample cannot have memorized a wrong label for it, so we check each given label with such a model. We shuffle $D$ and split it into $K$ cross-validation folds~\citep{stone1974crossvalidatory}, stratified on the given labels, and write $D_{k(i)}$ for the fold that contains sample $i$. We fit the base model, standardization included, on the other folds. The class it predicts for sample $i$ is the out-of-fold prediction $\hat{y}_i = f_{D \setminus D_{k(i)}}(x_i)$.

\paragraph{Keep rule and refit.}
We keep a sample only if its out-of-fold prediction equals its given label. The kept set is
\begin{equation}
  D_{\mathrm{keep}} = \{(x_i,\tilde{y}_i) \in D : \hat{y}_i = \tilde{y}_i\}.
  \label{eq:keep}
\end{equation}
We fit a new standardization and logistic regression on this set, which gives the refit model $f_{D_{\mathrm{keep}}}$ (Fig.~\ref{fig:teaser}a).

\paragraph{What the filter does not need.}
The filter needs no clean validation set or noise-rate estimate. It makes one pass over hard predictions, with no confidence threshold or per-class rule. The number of drops is the number of disagreements. We fix $K$ and the inverse regularization strength $C$ without tuning (Section~\ref{sec:setup}).
